%!TEX program = lualatex
% Relatório: Experimento 2 — Associação de Resistores

\documentclass[a4paper,12pt]{article}

% =========================================================
% Idioma e tipografia
% =========================================================
\usepackage[provide=*,brazilian]{babel}
\usepackage{iftex}
\ifLuaTeX
  \usepackage{fontspec}
  \setmainfont{Latin Modern Roman}
\else
  \usepackage[T1]{fontenc}
  \usepackage[utf8]{inputenc}
\fi

% =========================================================
% Layout
% =========================================================
\usepackage[top=3cm,bottom=2.5cm,left=3cm,right=2.5cm]{geometry}
\usepackage{setspace}
\onehalfspacing
\usepackage{indentfirst}
\setlength{\parindent}{1.25cm}

% =========================================================
% Matemática, tabelas, listas e figuras
% =========================================================
\usepackage{amsmath,amssymb}
\usepackage{float}
\usepackage{booktabs}
\usepackage{multirow}
\usepackage{array}
\usepackage{enumitem}
\usepackage{graphicx}

\setlist[itemize]{noitemsep,topsep=3pt}
\setlist[enumerate]{noitemsep,topsep=3pt}

\graphicspath{{./}{./Figs/}{./imagens/}{./figuras/}}

% Mantém o documento compilável mesmo quando alguma imagem ainda não está na pasta.
\newcommand{\imagemrelatorio}[2][width=.5\linewidth]{%
  \IfFileExists{#2}{%
    \includegraphics[#1]{#2}%
  }{%
    \fbox{\parbox[c][3.5cm][c]{.78\linewidth}{%
      \centering Imagem não encontrada: \texttt{\detokenize{#2}}%
    }}%
  }%
}

% =========================================================
% Links
% =========================================================
\usepackage{xurl}
\usepackage{hyperref}
\hypersetup{
  colorlinks=true,
  linkcolor=black,
  urlcolor=blue,
  citecolor=black,
  pdfauthor={José Mauro Elsas},
  pdftitle={Experimento 2 — Associação de Resistores}
}

% =========================================================
% Dados da capa
% =========================================================
\newcommand{\Instituicao}{Universidade do Estado do Rio de Janeiro}
\newcommand{\Curso}{Física VII -- Oceanografia}
\newcommand{\Professor}{Professor: Alexandre Magnus Gomes Carvalho}
\newcommand{\Aluno}{Alunos: Ricardo Renan,\\José Mauro Elsas}
\newcommand{\Cidade}{Rio de Janeiro}
\newcommand{\Ano}{2026}

\begin{document}
\pagenumbering{gobble}

% =========================================================
% Capa
% =========================================================
\begin{titlepage}
  \centering
  {\large \Instituicao\par}
  \vspace{0.5cm}
  {\large \Curso\par}
  \vspace{1cm}

  {\LARGE\textbf{Experimento 3}\par}
  \vspace{0.6cm}
  {\Large Associação de Resistores\par}

  \vspace{1cm}
  \imagemrelatorio[width=.3\linewidth]{uerj.jpg}

  \vfill

  \begin{flushright}
    \Aluno\\
    \Professor
  \end{flushright}

  \vspace{1cm}
  {\Cidade{} --- \Ano\par}
\end{titlepage}
\clearpage
\pagenumbering{arabic}

% =========================================================
% Objetivos
% =========================================================
\section{Objetivos}

São os seguintes os objetivos deste relatório:

\begin{enumerate}[label=\alph*.]
  \item Medir a resistência equivalente de pares de resistores associados em
  série e em paralelo;

  \item Comparar os valores experimentais obtidos com os respectivos valores
  teóricos.
\end{enumerate}

% =========================================================
% Introdução
% =========================================================
\section{Introdução}

O estudo quantitativo das associações de resistências desenvolveu-se no
século XIX, juntamente com a consolidação da teoria dos circuitos elétricos.
Em 1827, Georg Simon Ohm publicou \textit{Die galvanische Kette,
mathematisch bearbeitet}, obra na qual estabeleceu uma relação matemática
entre corrente elétrica, diferença de potencial e resistência, fornecendo uma
base para a análise dos circuitos resistivos \cite{ohm1827}. Posteriormente,
em 1845, Gustav Kirchhoff ampliou essa abordagem ao formular relações para
correntes e diferenças de potencial em redes elétricas constituídas por nós
e malhas, permitindo tratar sistematicamente circuitos mais complexos
\cite{kirchhoff_hist}. As associações em série e em paralelo podem, assim,
ser entendidas como casos fundamentais dessa análise de circuitos, nos quais
as leis de Ohm e de Kirchhoff conduzem diretamente às expressões utilizadas
para determinar a resistência equivalente.

A resistência equivalente de um conjunto de resistores depende dos valores
individuais dos componentes e da forma como eles são associados. Entre as
configurações mais simples estão as associações em série e em paralelo
\cite{openstax_resistores}.

Em uma associação em série, a corrente elétrica percorre sucessivamente todos
os resistores e apresenta o mesmo valor em cada componente. A resistência
equivalente é igual à soma das resistências individuais:

\[
  R_{\mathrm{eq,s}}=R_1+R_2+\cdots+R_n.
\]

Para dois resistores, essa expressão se reduz a

\[
  R_{\mathrm{eq,s}}=R_1+R_2.
\]

Consequentemente, a resistência equivalente de uma associação em série é maior
do que qualquer uma das resistências que a compõem.

Na associação em paralelo, os resistores estão submetidos à mesma diferença de
potencial, enquanto a corrente total se divide entre os diferentes ramos. A
resistência equivalente é determinada por

\[
  \frac{1}{R_{\mathrm{eq,p}}}
  =
  \frac{1}{R_1}
  +
  \frac{1}{R_2}
  +\cdots+
  \frac{1}{R_n}.
\]

Para dois resistores, pode-se escrever

\[
  R_{\mathrm{eq,p}}
  =
  \frac{R_1R_2}{R_1+R_2}.
\]

Nesse caso, a resistência equivalente é sempre menor do que a menor resistência
presente na associação \cite{openstax_resistores}.

Os resistores utilizados no experimento possuem tolerância nominal de
$\pm5\%$. Essa tolerância indica que o valor real de cada componente pode
diferir do respectivo valor nominal dentro desse intervalo. Como todos os
resistores empregados possuem a mesma tolerância relativa, os limites da
resistência equivalente podem ser obtidos aplicando-se os fatores $0{,}95$ e
$1{,}05$ ao valor equivalente calculado a partir das resistências nominais:

\[
  R_{\min}=0{,}95R_{\mathrm{teo}}
  \qquad\text{e}\qquad
  R_{\max}=1{,}05R_{\mathrm{teo}}.
\]

Dessa forma, a comparação entre os valores calculados e medidos permite
verificar experimentalmente as relações para associações em série e em paralelo
e avaliar se as diferenças observadas são compatíveis com a tolerância dos
componentes.

% =========================================================
% Materiais e métodos
% =========================================================
\section{Materiais e Métodos}

Neste experimento foram utilizados os seguintes materiais:

\begin{enumerate}[label=\alph*.]
  \item Multímetro digital Minipa ET-1110A e cabos de conexão.

  \begin{figure}[H]
    \centering
    \imagemrelatorio[width=.72\linewidth]{Figs/multimetro_resistores.jpeg}
    \caption{Multímetro digital e cabos utilizados nas medições.}
    \label{fig:multimetro}
  \end{figure}

  \item \textit{Protoboard} com bornes para montagem das associações.

  \begin{figure}[H]
    \centering
    \imagemrelatorio[width=.65\linewidth]{Figs/placa_associacao.jpeg}
    \caption{\textit{Protoboard} utilizada na montagem das associações.}
    \label{fig:placa_associacao}
  \end{figure}

  \item Resistores de valores nominais
  $47\,\Omega$, $100\,\Omega$, $150\,\Omega$, $330\,\Omega$ e
  $470\,\Omega$, todos com tolerância nominal de $\pm5\%$, além dos
  cabos e conectores empregados nas montagens.

  \begin{figure}[H]
    \centering
    \imagemrelatorio[height = 12 cm]{Figs/resistores_associacoes.jpeg}
    \caption{Resistores e conexões utilizados no experimento.}
    \label{fig:resistores_associacoes}
  \end{figure}
\end{enumerate}

\subsection{Procedimento experimental}

Manteve-se fixo um resistor de valor nominal

\[
  R_1=47\,\Omega,
\]

associando-o sucessivamente aos resistores

\[
  R_2=100\,\Omega,\quad
  150\,\Omega,\quad
  330\,\Omega
  \quad\text{e}\quad
  470\,\Omega.
\]

Na primeira etapa, cada par foi montado em série e a resistência equivalente
foi medida diretamente com o multímetro na função de ohmímetro. Os valores
experimentais foram comparados com os valores teóricos calculados por

\[
  R_{\mathrm{teo,s}}=R_1+R_2.
\]

Na segunda etapa, os mesmos pares foram montados em paralelo. As resistências
equivalentes foram novamente medidas com o multímetro e comparadas com os
valores calculados por

\[
  R_{\mathrm{teo,p}}
  =
  \frac{R_1R_2}{R_1+R_2}.
\]

Para quantificar a diferença entre cada valor medido e o respectivo valor
teórico, utilizou-se

\[
  \varepsilon_{\%}
  =
  \frac{\lvert R_{\mathrm{exp}}-R_{\mathrm{teo}}\rvert}
       {R_{\mathrm{teo}}}
  \times100.
\]

Não foram calculados média ou desvio-padrão, pois foi registrada uma medição
da resistência equivalente para cada configuração. Também não foram utilizados
gráficos, uma vez que a comparação direta entre valores teóricos, experimentais
e faixas de tolerância é suficientemente representada pelas tabelas.

% =========================================================
% Resultados
% =========================================================
\section{Resultados}

\subsection{Associação em série}

Para a associação em série, aplicou-se

\[
  R_{\mathrm{teo,s}}=R_1+R_2.
\]

As resistências equivalentes teóricas foram calculadas individualmente:

\[
\begin{aligned}
R_{\mathrm{teo,s}}(47,100)
  &=47+100
   =147\,\Omega,\\[0.2cm]
R_{\mathrm{teo,s}}(47,150)
  &=47+150
   =197\,\Omega,\\[0.2cm]
R_{\mathrm{teo,s}}(47,330)
  &=47+330
   =377\,\Omega,\\[0.2cm]
R_{\mathrm{teo,s}}(47,470)
  &=47+470
   =517\,\Omega.
\end{aligned}
\]

Os valores experimentais obtidos foram, respectivamente,
$146{,}8\,\Omega$, $197{,}6\,\Omega$, $391\,\Omega$ e
$524\,\Omega$.

As diferenças percentuais entre os valores experimentais e teóricos foram:

\[
\begin{aligned}
\varepsilon_{47+100}
  &=
  \frac{|146{,}8-147|}{147}\times100
  =0{,}14\%,\\[0.2cm]
\varepsilon_{47+150}
  &=
  \frac{|197{,}6-197|}{197}\times100
  =0{,}30\%,\\[0.2cm]
\varepsilon_{47+330}
  &=
  \frac{|391-377|}{377}\times100
  =3{,}71\%,\\[0.2cm]
\varepsilon_{47+470}
  &=
  \frac{|524-517|}{517}\times100
  =1{,}35\%.
\end{aligned}
\]

Para comparação com a tolerância dos componentes, utilizaram-se os limites

\[
  R_{\min}=0{,}95R_{\mathrm{teo}}
  \qquad\text{e}\qquad
  R_{\max}=1{,}05R_{\mathrm{teo}}.
\]

Assim, as faixas teóricas de tolerância são:

\[
\begin{aligned}
147\,\Omega &:\
139{,}65\,\Omega
\leq R \leq
154{,}35\,\Omega,\\
197\,\Omega &:\
187{,}15\,\Omega
\leq R \leq
206{,}85\,\Omega,\\
377\,\Omega &:\
358{,}15\,\Omega
\leq R \leq
395{,}85\,\Omega,\\
517\,\Omega &:\
491{,}15\,\Omega
\leq R \leq
542{,}85\,\Omega.
\end{aligned}
\]

\begin{table}[H]
  \centering
  \caption{Resultados para as associações de resistores em série.}
  \label{tab:serie}
  \medskip

  \footnotesize
  \setlength{\tabcolsep}{4pt}
  \begin{tabular}{ccccc}
    \toprule
    \shortstack{$R_2$\\($\Omega$)}
    & \shortstack{$R_{\mathrm{teo}}$\\($\Omega$)}
    & \shortstack{$R_{\mathrm{exp}}$\\($\Omega$)}
    & \shortstack{\textbf{Diferença}\\\textbf{percentual}}
    & \shortstack{\textbf{Faixa teórica}\\\textbf{com tolerância de 5\%}} \\
    \midrule
    100 & 147 & 146,8 & 0,14\% & 139,65 a 154,35 \\
    150 & 197 & 197,6 & 0,30\% & 187,15 a 206,85 \\
    330 & 377 & 391   & 3,71\% & 358,15 a 395,85 \\
    470 & 517 & 524   & 1,35\% & 491,15 a 542,85 \\
    \bottomrule
  \end{tabular}
\end{table}

Todos os valores experimentais obtidos para as associações em série ficaram
dentro das respectivas faixas correspondentes à tolerância de $\pm5\%$. A maior
diferença percentual ocorreu na associação de $47\,\Omega$ com $330\,\Omega$,
para a qual foram calculados $377\,\Omega$ e medidos $391\,\Omega$, resultando
em diferença de $3{,}71\%$.

\subsection{Associação em paralelo}

Para dois resistores associados em paralelo, empregou-se

\[
  R_{\mathrm{teo,p}}
  =
  \frac{R_1R_2}{R_1+R_2}.
\]

Para o par de $47\,\Omega$ e $100\,\Omega$:

\[
\begin{aligned}
R_{\mathrm{teo,p}}
  &=
  \frac{47\times100}{47+100}\\
  &=
  \frac{4700}{147}\\
  &=31{,}97\,\Omega.
\end{aligned}
\]

Para o par de $47\,\Omega$ e $150\,\Omega$:

\[
\begin{aligned}
R_{\mathrm{teo,p}}
  &=
  \frac{47\times150}{47+150}\\
  &=
  \frac{7050}{197}\\
  &=35{,}79\,\Omega.
\end{aligned}
\]

Para o par de $47\,\Omega$ e $330\,\Omega$:

\[
\begin{aligned}
R_{\mathrm{teo,p}}
  &=
  \frac{47\times330}{47+330}\\
  &=
  \frac{15510}{377}\\
  &=41{,}14\,\Omega.
\end{aligned}
\]

Para o par de $47\,\Omega$ e $470\,\Omega$:

\[
\begin{aligned}
R_{\mathrm{teo,p}}
  &=
  \frac{47\times470}{47+470}\\
  &=
  \frac{22090}{517}\\
  &=42{,}73\,\Omega.
\end{aligned}
\]

Os valores experimentais foram, respectivamente,
$32{,}2\,\Omega$, $36{,}4\,\Omega$, $41{,}7\,\Omega$ e
$43{,}3\,\Omega$.

As diferenças percentuais foram calculadas por:

\[
\begin{aligned}
\varepsilon_{47\parallel100}
  &=
  \frac{|32{,}2-31{,}97|}{31{,}97}\times100
  \approx0{,}71\%,\\[0.2cm]
\varepsilon_{47\parallel150}
  &=
  \frac{|36{,}4-35{,}79|}{35{,}79}\times100
  \approx1{,}71\%,\\[0.2cm]
\varepsilon_{47\parallel330}
  &=
  \frac{|41{,}7-41{,}14|}{41{,}14}\times100
  \approx1{,}36\%,\\[0.2cm]
\varepsilon_{47\parallel470}
  &=
  \frac{|43{,}3-42{,}73|}{42{,}73}\times100
  \approx1{,}34\%.
\end{aligned}
\]

As faixas correspondentes à tolerância de $\pm5\%$ são:

\[
\begin{aligned}
31{,}97\,\Omega &:\
30{,}37\,\Omega
\leq R \leq
33{,}57\,\Omega,\\
35{,}79\,\Omega &:\
34{,}00\,\Omega
\leq R \leq
37{,}58\,\Omega,\\
41{,}14\,\Omega &:\
39{,}08\,\Omega
\leq R \leq
43{,}20\,\Omega,\\
42{,}73\,\Omega &:\
40{,}59\,\Omega
\leq R \leq
44{,}86\,\Omega.
\end{aligned}
\]

\begin{table}[H]
  \centering
  \caption{Resultados para as associações de resistores em paralelo.}
  \label{tab:paralelo}
  \medskip

  \footnotesize
  \setlength{\tabcolsep}{4pt}
  \begin{tabular}{ccccc}
    \toprule
    \shortstack{$R_2$\\($\Omega$)}
    & \shortstack{$R_{\mathrm{teo}}$\\($\Omega$)}
    & \shortstack{$R_{\mathrm{exp}}$\\($\Omega$)}
    & \shortstack{\textbf{Diferença}\\\textbf{percentual}}
    & \shortstack{\textbf{Faixa teórica}\\\textbf{com tolerância de 5\%}} \\
    \midrule
    100 & 31,97 & 32,2 & 0,71\% & 30,37 a 33,57 \\
    150 & 35,79 & 36,4 & 1,71\% & 34,00 a 37,58 \\
    330 & 41,14 & 41,7 & 1,36\% & 39,08 a 43,20 \\
    470 & 42,73 & 43,3 & 1,34\% & 40,59 a 44,86 \\
    \bottomrule
  \end{tabular}
\end{table}

Também nas associações em paralelo todos os valores medidos permaneceram
dentro das respectivas faixas de tolerância. Além disso, em todas as
configurações a resistência equivalente experimental foi inferior a
$47\,\Omega$, que era o menor valor de resistência individual em cada par,
em concordância com o comportamento previsto para resistores em paralelo.

Comparando-se as duas formas de associação, os resultados em série apresentaram
resistências equivalentes maiores do que as resistências individuais, enquanto
os resultados em paralelo apresentaram resistências equivalentes menores que
$47\,\Omega$. Esse comportamento constitui uma verificação adicional da
coerência das montagens e das medições realizadas.

% =========================================================
% Conclusão
% =========================================================
\section{Conclusão}

O experimento permitiu medir a resistência equivalente de quatro pares de
resistores montados em série e em paralelo e comparar os resultados com os
valores obtidos pelas relações teóricas correspondentes.

Nas associações em série, foram calculadas resistências equivalentes de
$147\,\Omega$, $197\,\Omega$, $377\,\Omega$ e $517\,\Omega$, enquanto os
valores medidos foram $146{,}8\,\Omega$, $197{,}6\,\Omega$, $391\,\Omega$ e
$524\,\Omega$. As diferenças percentuais foram, respectivamente,
$0{,}14\%$, $0{,}30\%$, $3{,}71\%$ e $1{,}35\%$.

Nas associações em paralelo, os valores teóricos obtidos foram
$31{,}97\,\Omega$, $35{,}79\,\Omega$, $41{,}14\,\Omega$ e
$42{,}73\,\Omega$, comparados aos valores experimentais de
$32{,}2\,\Omega$, $36{,}4\,\Omega$, $41{,}7\,\Omega$ e
$43{,}3\,\Omega$. As diferenças percentuais ficaram entre
$0{,}71\%$ e $1{,}71\%$.

Todos os valores experimentais ficaram dentro das faixas correspondentes à
tolerância nominal de $\pm5\%$ dos resistores. A maior diferença observada foi
de $3{,}71\%$, na associação em série dos resistores de $47\,\Omega$ e
$330\,\Omega$, permanecendo ainda assim dentro da faixa esperada.

Os resultados também reproduziram as características fundamentais das duas
associações: em série, a resistência equivalente foi maior do que cada
resistência individual, enquanto em paralelo à resistência equivalente foi
menor do que o menor resistor de cada par. Conclui-se, portanto, que as
medições realizadas foram compatíveis com os valores teóricos e com as
tolerâncias especificadas para os componentes utilizados.

\newpage

% =========================================================
% Referências
% =========================================================
\begin{thebibliography}{99}

\bibitem{openstax_resistores}
LING, Samuel J.; MOEBS, William; SANNY, Jeff.
\textit{University Physics Volume 2}.
Houston: OpenStax, 2016.
Seção 10.2: Resistors in Series and Parallel.
Disponível em:
\url{https://openstax.org/books/university-physics-volume-2/pages/10-2-resistors-in-series-and-parallel}.
Acesso em: 22 ago. 2026.
% =========================================================
% Referências
% =========================================================


\bibitem{ohm1827}
OHM, Georg Simon.
\textit{Die galvanische Kette, mathematisch bearbeitet}.
Berlin: T. H. Riemann, 1827.
Registro disponível em: Smithsonian Libraries and Archives.
Disponível em:
\url{https://www.si.edu/object/die-galvanische-kette-mathematisch-bearbeitet-von-dr-gs-ohm:siris_sil_227673}.
Acesso em: 22 ago. 2026.

\bibitem{kirchhoff_hist}
O'CONNOR, J. J.; ROBERTSON, E. F.
\textit{Gustav Robert Kirchhoff}.
MacTutor History of Mathematics Archive.
University of St Andrews.
Disponível em:
\url{https://mathshistory.st-andrews.ac.uk/Biographies/Kirchhoff/}.
Acesso em: 22 ago. 2026.

\bibitem{openstax_resistores}
LING, Samuel J.; MOEBS, William; SANNY, Jeff.
\textit{University Physics Volume 2}.
Houston: OpenStax, 2016.
Seção 10.2: Resistors in Series and Parallel.
Disponível em:
\url{https://openstax.org/books/university-physics-volume-2/pages/10-2-resistors-in-series-and-parallel}.
Acesso em: 22 ago. 2026.
\end{thebibliography}

% =========================================================
% Anexo
% =========================================================
\section{Anexo}

\begin{figure}[H]
  \centering
  \imagemrelatorio[height=0.78\textheight,keepaspectratio]{Figs/dados_associacoes.jpeg}
  \caption{Dados obtidos no laboratório.}
  \label{fig:dados_laboratorio}
\end{figure}

\end{document}

